% =====================================================================
% 1. METHODOLOGY TO DEVELOP MACHINE LEARNING MODELS
% =====================================================================
\section{Methodology to develop Machine Learning models}

\subsection{Methodological guidelines for developing Machine Learning models}

The project predicts the grade of TCGA gliomas (LGG vs GBM) and the age at diagnosis from clinical
data and the mutation status of 20 genes, and explores the mutation profiles without labels (the
theory and the critical analysis are in the companion PRDL report). The development followed five
guidelines, each enforced by the code rather than by convention:

\begin{enumerate}
  \item \textbf{Staged, reproducible workflow.} Five notebooks run in order (data preparation, EDA,
        unsupervised analysis, regression, classification). Each one reads only the output of the
        previous stages, fixes the seed (42) for every random process, and writes its results to
        \texttt{results/}: figures, \LaTeX{} tables and a JSON summary per stage. Both reports are built
        from these files: every number in the text is a macro generated from the JSON summaries by
        \texttt{report/\allowbreak{}make\_\allowbreak{}numbers.\allowbreak{}py}, so no number is typed by hand.
  \item \textbf{Leakage prevention by construction.} Columns that encode the target (TCGA project,
        histological diagnosis) are excluded from the predictor lists in \texttt{src/data.py}, and the
        classification notebook stops if one of them appears. All preprocessing that learns from the
        data (standardisation, one-hot encoding, discretisation) is a step of a scikit-learn
        \texttt{Pipeline}, so it is refitted on every training fold \cite{kaufman2012}.
  \item \textbf{Nested cross-validation for every supervised model} (scheme of the course notebook
        \texttt{Example\_\allowbreak{}ExtrOSA\_\allowbreak{}ModelTuning\_\allowbreak{}NestedLoop}): 10 outer folds for evaluation and 3 inner
        folds for \texttt{GridSearchCV}, stratified for classification
        (Figure~\ref{fig:mllb-nested}).
  \item \textbf{Evaluation against a baseline, with uncertainty.} Out-of-fold predictions for every
        patient; bootstrap 95\% intervals; no model is declared better than another when the intervals
        overlap, and a paired bootstrap of the difference complements the comparison.
  \item \textbf{Verification of AI-generated code.} Every notebook was executed in a clean copy of
        the repository before delivery, every figure was inspected, and the run on the author's
        computer was compared value by value with that run (Appendix A).
\end{enumerate}

\begin{figure}[H]
\centering
\begin{tabular}{|p{0.92\textwidth}|}
\hline
\textbf{for} each outer fold $k = 1,\dots,10$ (stratified, seed 42):\\
\quad split patients into training part $T_k$ (90\%) and test fold $V_k$ (10\%)\\
\quad \textbf{for} each hyperparameter configuration $\theta$ in the grid:\\
\qquad 3-fold CV inside $T_k$: fit pipeline (preprocessing + model) on 2/3, score on 1/3\\
\quad choose $\theta^*_k$ with the best mean inner score (AUC or RMSE)\\
\quad refit the pipeline with $\theta^*_k$ on all of $T_k$\\
\quad predict $V_k$ $\rightarrow$ out-of-fold predictions; permutation importance on $V_k$\\
metrics and bootstrap intervals on the \NClean{} out-of-fold predictions\\
final model for interpretation: the same grid search applied once to all patients\\
\hline
\end{tabular}
\caption{Nested cross-validation as implemented in \texttt{src/\allowbreak{}evaluation.\allowbreak{}py} (\texttt{nested\_cv},
\texttt{fit\_final}).}
\label{fig:mllb-nested}
\end{figure}

\subsection{Tools to develop Machine Learning models}

Table~\ref{tab:mllb-tools} lists the software. The environment is created with Anaconda from
\texttt{environment.yml} (Python 3.11); all libraries are pinned in \texttt{requirements.\allowbreak{}txt} with
exact versions that were installed and tested together; the script \texttt{check\_env.py} verifies them
and writes the result to the
\texttt{results/logs/} folder. Notebooks were run in Visual Studio Code with a
dedicated Jupyter kernel; the code is versioned with git and published on GitHub
(Section~\ref{sec:mllb-code}).

\begin{table}[H]
\centering
\caption{Software used. Versions are those pinned in \texttt{requirements.\allowbreak{}txt}.}
\label{tab:mllb-tools}
\begin{tabularx}{\textwidth}{l l X}
\toprule
\textbf{Tool} & \textbf{Version} & \textbf{Use} \\
\midrule
Python & 3.11 & Language \\
pandas \cite{mckinney2010} & 2.2.3 & Data loading, cleaning and tables \\
NumPy \cite{harris2020} & 2.2.6 & Numerical arrays, bootstrap \\
SciPy \cite{virtanen2020} & 1.15.3 & Statistical tests, hierarchical clustering \\
scikit-learn \cite{pedregosa2011} & 1.6.1 & Pipelines, models, cross-validation, PCA, MDS, t-SNE, $k$-means \\
XGBoost \cite{chen2016} & 2.1.4 & Gradient boosting \\
statsmodels \cite{seabold2010} & 0.14.4 & OLS inference, VIF, influence measures, multiple testing \\
SHAP \cite{lundberg2017} & 0.48.0 & SHAP values of the winning models (TreeSHAP, linear) \\
Matplotlib \cite{hunter2007} / seaborn \cite{waskom2021} & 3.10.3 / 0.13.2 & Figures \\
Jupyter (ipykernel, nbconvert) & 6.29.5 / 7.16.6 & Notebooks and their execution \\
Anaconda, VS Code, git/GitHub & -- & Environment, editor, version control \\
Claude (Anthropic) & -- & AI code assistant (Appendix A) \\
\bottomrule
\end{tabularx}
\end{table}

The code is organised as notebooks that call functions from a small package, so that the same code
is used everywhere and can be read on its own during the defense:
\texttt{config.py} (paths, seed, validation settings), \texttt{data.py} (cleaning, encoding, feature
sets), \texttt{eda.py} (tests and effect sizes), \texttt{unsupervised.py}, \texttt{pipelines.py}
(preprocessing), \texttt{models.py} (pipelines and grids), \texttt{evaluation.py} (nested CV, metrics,
bootstrap, comparisons), \texttt{explain.py} (SHAP values and plots), \texttt{causal.py} (mediation
decomposition, robust OLS, causal diagram), \texttt{plotting.py} (style and fixed colours: LGG blue, GBM orange in every
figure) and \texttt{reporting.py} (\LaTeX{} tables and JSON summaries).
